\part{\trFukayaCategory}
\markboth{II　\trFukayaCategory}{II　\trFukayaCategory}
本章的主要目的，是借助由\trGromov{}与\trFloer{}所开创的经典分析方法，在\trExactSymplectic\trManifold{}这个相对简单的特殊情形下，讲解\trFukaya{}的 \trACategory{}的定义。可以说，这条进路如今已经过时了：\trFukaya--\trOh--\trOhta--\trOno{}的基础性工作 [60]、[56] 已取而代之，它使人能够在完全一般的情形下定义\trFukayaCategory。然而，这个一般的定义无论在分析上还是代数上都相当繁复，而其中的许多精巧之处，在我们这里所讨论的情形下本来是用不上的。

为了照顾非专业读者，我们先从\trExact\trLagrange\trSubmanifold{}的\trFloerCohomology{}群的原始定义 [44] 讲起。人们自 20 世纪 90 年代初起就已经知道，\trFloer{}群容许一些由一类适当的\trPointedBoundary{}的\trRiemann\trSurface{}所标号的运算，这使它们成为一个\trOpenString{} TQFT\footnote{译注：TQFT，即 topological quantum field theory（\trTopology{}量子场论）。}（这一想法似乎源自\trDonaldson；处理它或其\trClosedString{}类似物的论文有 [119]、[106]、[135]、[131]）。\trFukaya{} \trACategory{}首见于 [52]，它出现在容许\trComplexStructure{}可变的\trRiemann\trSurface{}族之时（实际上它只用到\trPointedBoundary{}的\trDisk，因此应当被看作某个更复杂理论的第一块构件，而这个理论在很大程度上尚待开发）。所谓经典进路，指的是用一个零阶\trInhomogeneous{}项去\trPerturbation{}\trPseudoholomorphic{}映射方程，使解的\trModuliSpace{}变得\trSmooth（从而避免使用\trVirtualFundamentalChain）。与 [116] 一样，\trInhomogeneous{}项随\trRiemann\trSurface{}的\trComplexStructure{}而变，而我们会花一些篇幅解释其确切的要求是什么。一旦把这一点澄清，那些基本的分析结果便没有什么特别困难或新颖之处，因此我们陈述它们时只配以最简的证明概要。相比之下，我们会把更多的注意力放在\trGrading{}与正负号的问题上，因为这两者对我们的应用都很重要。关于\trGrading，我们采用\trKontsevich{}的\trGraded\trLagrange\trSubmanifold{}形式体系 [90]、[125]，它比\trFloer{}与\trViterbo{}较老的\trRelative{}\trGrading{}多保留一点信息。研究\trFloer{}理论中正负号问题的一般框架在 [48] 中建立，\trLagrange{}\trFloerCohomology{}的情形则在 [60] 中得到处理；但我们仍然觉得值得把该理论的一部分重新做一遍。

这里处理的论题还有：带\trInvolution{}的\trSymplectic\trManifold{}的\trEquivariant{}\trFukayaCategory，以及一个\trSurface{}的\trFukayaCategory{}的（部分）组合性质。后一论题已经在若干语境中出现过：\trElliptic{}曲线的\trHomology\trMirrorSymmetry{} [90]、[108]、[107]，\trSurface{}上曲线的\trFloerHomology{} [135]、[35]，以及\trLegendrian{}纽结理论 [30]、[43]。我们将集中讨论把分析方法与组合方法联系起来的\trRegularity{}论证，因为在文献中这一方面一直较少被顾及。
